\documentclass[tikz,border=3mm,multi=tikzpicture]{standalone}
\usepackage{eleve-geometria}

\begin{document}

% FIGURA 1 — fig-01-elementos-da-circunferencia
\begin{tikzpicture}[eleve figura, x=1cm, y=1cm]
  \path[use as bounding box] (-0.15,-0.10) rectangle (7.75,7.35);
  \coordinate (O) at (3.80,3.85);
  \def\R{2.55}
  \fill[eleveazulclaro] (O) circle (\R);
  \draw[eleve aresta, line width=1.7pt] (O) circle (\R);
  \fill[eleveazul] (O) circle (2.1pt);
  \node[font=\large, text=elevecinza] at ($(O)+(-.30,-.32)$) {$O$};

  % Diâmetro e raio
  \coordinate (D1) at ($(O)+(-\R,0)$);
  \coordinate (D2) at ($(O)+(\R,0)$);
  \draw[eleve destaque, line width=1.7pt] (D1) -- (D2);
  \node[font=\large\bfseries, text=elevelaranja] at (2.45,4.55) {diâmetro};
  \coordinate (R) at ($(O)+(58:\R)$);
  \draw[eleve aresta, line width=1.6pt] (O) -- (R);
  \node[font=\large\bfseries, text=eleveazul] at (5.18,5.20) {raio};

  % Corda e flecha
  \coordinate (A) at ($(O)+(215:\R)$);
  \coordinate (B) at ($(O)+(325:\R)$);
  \coordinate (M) at ($(A)!0.5!(B)$);
  \coordinate (F) at ($(O)+(270:\R)$);
  \draw[eleve aresta, line width=1.5pt] (A) -- (B);
  \node[font=\large\bfseries, text=eleveazul] at ($(M)+(0,0.82)$) {corda};
  \draw[eleve destaque, |-|, line width=1.5pt] (M) -- (F);
  \node[font=\large\bfseries, text=elevelaranja, right=5pt] at ($(M)!0.5!(F)$)
    {flecha};

  % Arco destacado numa região livre
  \draw[eleve destaque, line width=2.8pt]
    ($(O)+(112:\R)$) arc[start angle=112,end angle=166,radius=\R];
  \node[font=\large\bfseries, text=elevelaranja] at (1.10,6.47) {arco};
  \draw[eleve auxiliar, ->] (1.68,6.32) -- (2.15,5.92);
\end{tikzpicture}

% FIGURA 2 — fig-02-posicoes-de-um-ponto
\begin{tikzpicture}[eleve figura, x=1cm, y=1cm]
  \path[use as bounding box] (-0.15,-0.05) rectangle (7.75,10.85);

  % Ponto interno
  \coordinate (O1) at (3.10,9.00);
  \fill[eleveazulclaro] (O1) circle (1.35);
  \draw[eleve aresta, line width=1.5pt] (O1) circle (1.35);
  \fill[eleveazul] (O1) circle (2pt) node[font=\large, below left=1pt] {$O$};
  \coordinate (P1) at ($(O1)+(25:0.76)$);
  \draw[eleve auxiliar] (O1) -- (P1);
  \fill[elevelaranja] (P1) circle (2.5pt)
    node[font=\large\bfseries, text=elevelaranja, above right=2pt] {$P$};
  \node[font=\large\bfseries, text=elevecinza, align=left] at (6.05,9.00)
    {posição\\interna};

  % Ponto pertencente
  \coordinate (O2) at (3.10,5.45);
  \fill[eleveazulclaro] (O2) circle (1.35);
  \draw[eleve aresta, line width=1.5pt] (O2) circle (1.35);
  \fill[eleveazul] (O2) circle (2pt) node[font=\large, below left=1pt] {$O$};
  \coordinate (P2) at ($(O2)+(28:1.35)$);
  \draw[eleve auxiliar] (O2) -- (P2);
  \fill[elevelaranja] (P2) circle (2.5pt)
    node[font=\large\bfseries, text=elevelaranja, above right=2pt] {$P$};
  \node[font=\large\bfseries, text=elevecinza, align=left] at (6.05,5.45)
    {na\\circunferência};

  % Ponto externo
  \coordinate (O3) at (3.10,1.90);
  \fill[eleveazulclaro] (O3) circle (1.35);
  \draw[eleve aresta, line width=1.5pt] (O3) circle (1.35);
  \fill[eleveazul] (O3) circle (2pt) node[font=\large, below left=1pt] {$O$};
  \coordinate (P3) at ($(O3)+(24:2.08)$);
  \draw[eleve auxiliar] (O3) -- (P3);
  \fill[elevelaranja] (P3) circle (2.5pt)
    node[font=\large\bfseries, text=elevelaranja, above right=2pt] {$P$};
  \node[font=\large\bfseries, text=elevecinza, align=left] at (6.05,1.90)
    {posição\\externa};
\end{tikzpicture}

% FIGURA 3 — fig-03-posicoes-de-uma-reta
\begin{tikzpicture}[eleve figura, x=1cm, y=1cm]
  \path[use as bounding box] (-0.15,-0.05) rectangle (8.75,11.15);

  % Externa
  \coordinate (O1) at (3.30,9.20);
  \fill[eleveazulclaro] (O1) circle (1.35);
  \draw[eleve aresta, line width=1.5pt] (O1) circle (1.35);
  \draw[eleve destaque, <->, line width=1.6pt] (0.85,10.86) -- (5.75,10.86);
  \draw[eleve auxiliar, <->] (O1) -- (3.30,10.86);
  \node[font=\large\bfseries, text=elevecinza, align=left] at (7.15,9.20)
    {externa\\nenhum ponto};

  % Tangente
  \coordinate (O2) at (3.30,5.55);
  \fill[eleveazulclaro] (O2) circle (1.35);
  \draw[eleve aresta, line width=1.5pt] (O2) circle (1.35);
  \coordinate (T) at (3.30,6.90);
  \draw[eleve destaque, <->, line width=1.6pt] (0.85,6.90) -- (5.75,6.90);
  \draw[eleve auxiliar] (O2) -- (T);
  \draw[elevecinza, line width=1pt] (3.30,6.65) -- (3.55,6.65) -- (3.55,6.90);
  \fill[elevelaranja] (T) circle (2.5pt);
  \node[font=\large\bfseries, text=elevecinza, align=left] at (7.15,5.55)
    {tangente\\um ponto};

  % Secante
  \coordinate (O3) at (3.30,1.90);
  \fill[eleveazulclaro] (O3) circle (1.35);
  \draw[eleve aresta, line width=1.5pt] (O3) circle (1.35);
  \draw[eleve destaque, <->, line width=1.6pt] (0.85,1.55) -- (5.75,2.25);
  \fill[elevelaranja] (1.99,1.71) circle (2.5pt);
  \fill[elevelaranja] (4.61,2.09) circle (2.5pt);
  \node[font=\large\bfseries, text=elevecinza, align=left] at (7.15,1.90)
    {secante\\dois pontos};
\end{tikzpicture}

% FIGURA 4 — fig-04-circunferencias-externas
\begin{tikzpicture}[eleve figura, x=1cm, y=1cm]
  \path[use as bounding box] (-0.15,-0.05) rectangle (9.15,10.75);

  % Externas
  \coordinate (O1) at (2.35,9.00);
  \coordinate (O2) at (6.55,9.00);
  \draw[eleve auxiliar] (O1) -- (O2);
  \fill[eleveazulclaro] (O1) circle (1.25);
  \fill[orange!16] (O2) circle (1.05);
  \draw[eleve aresta, line width=1.5pt] (O1) circle (1.25);
  \draw[eleve destaque, line width=1.5pt] (O2) circle (1.05);
  \fill[eleveazul] (O1) circle (2pt) node[font=\large, below=2pt] {$O_1$};
  \fill[elevelaranja] (O2) circle (2pt) node[font=\large, below=2pt] {$O_2$};
  \node[font=\large\bfseries, text=elevecinza] at (4.45,7.38) {externas};

  % Tangentes externas
  \coordinate (O3) at (2.55,5.25);
  \coordinate (O4) at (5.85,5.25);
  \coordinate (T) at (4.25,5.25);
  \draw[eleve auxiliar] (O3) -- (O4);
  \fill[eleveazulclaro] (O3) circle (1.70);
  \fill[orange!16] (O4) circle (1.60);
  \draw[eleve aresta, line width=1.5pt] (O3) circle (1.70);
  \draw[eleve destaque, line width=1.5pt] (O4) circle (1.60);
  \fill[eleveazul] (O3) circle (2pt) node[font=\large, below=2pt] {$O_1$};
  \fill[elevelaranja] (O4) circle (2pt) node[font=\large, below=2pt] {$O_2$};
  \fill[elevecinza] (T) circle (2.5pt) node[font=\large, above=3pt] {$T$};
  \node[font=\large\bfseries, text=elevecinza] at (4.20,3.25)
    {tangentes externas};

  % Secantes
  \coordinate (O5) at (3.20,1.35);
  \coordinate (O6) at (5.50,1.35);
  \fill[eleveazulclaro, opacity=.72] (O5) circle (1.35);
  \fill[orange!20, opacity=.72] (O6) circle (1.35);
  \draw[eleve aresta, line width=1.5pt] (O5) circle (1.35);
  \draw[eleve destaque, line width=1.5pt] (O6) circle (1.35);
  \fill[eleveazul] (O5) circle (2pt) node[font=\large, below=2pt] {$O_1$};
  \fill[elevelaranja] (O6) circle (2pt) node[font=\large, below=2pt] {$O_2$};
  \node[font=\large\bfseries, text=elevecinza] at (4.35,-0.02) {secantes};
\end{tikzpicture}

% FIGURA 5 — fig-05-circunferencias-internas
\begin{tikzpicture}[eleve figura, x=1cm, y=1cm]
  \path[use as bounding box] (-0.15,-0.85) rectangle (8.85,11.65);

  % Tangentes internas
  \coordinate (O1) at (4.20,9.65);
  \coordinate (O2) at (5.45,9.65);
  \fill[eleveazulclaro] (O1) circle (2.20);
  \fill[orange!18] (O2) circle (0.95);
  \draw[eleve aresta, line width=1.5pt] (O1) circle (2.20);
  \draw[eleve destaque, line width=1.5pt] (O2) circle (0.95);
  \draw[eleve auxiliar] (O1) -- (O2);
  \fill[eleveazul] (O1) circle (2pt) node[font=\large, below left=2pt] {$O_1$};
  \fill[elevelaranja] (O2) circle (2pt) node[font=\large, below left=2pt] {$O_2$};
  \fill[elevecinza] (6.40,9.65) circle (2.5pt)
    node[font=\large, right=3pt] {$T$};
  \node[font=\large\bfseries, text=elevecinza, align=left, anchor=west]
    at (6.72,9.65) {tangentes\\internas};

  % Internas
  \coordinate (O3) at (4.20,4.85);
  \coordinate (O4) at (5.10,5.25);
  \fill[eleveazulclaro] (O3) circle (2.05);
  \fill[orange!18] (O4) circle (0.80);
  \draw[eleve aresta, line width=1.5pt] (O3) circle (2.05);
  \draw[eleve destaque, line width=1.5pt] (O4) circle (0.80);
  \draw[eleve auxiliar] (O3) -- (O4);
  \fill[eleveazul] (O3) circle (2pt) node[font=\large, below left=2pt] {$O_1$};
  \fill[elevelaranja] (O4) circle (2pt) node[font=\large, above right=2pt] {$O_2$};
  \node[font=\large\bfseries, text=elevecinza, anchor=west]
    at (6.52,4.85) {internas};

  % Concêntricas
  \coordinate (O5) at (3.30,0.80);
  \fill[eleveazulclaro] (O5) circle (1.50);
  \fill[orange!14] (O5) circle (0.78);
  \draw[eleve aresta, line width=1.5pt] (O5) circle (1.50);
  \draw[eleve destaque, line width=1.5pt] (O5) circle (0.78);
  \fill[elevecinza] (O5) circle (2.3pt);
  \draw[eleve auxiliar, ->] (5.00,0.80) -- (3.42,0.80);
  \node[font=\large, text=elevecinza, align=left, anchor=west] at (5.20,0.80)
    {$O_1=O_2$\\\textbf{concêntricas}};
\end{tikzpicture}

% FIGURA 6 — fig-06-angulo-central-e-inscrito
\begin{tikzpicture}[eleve figura, x=1cm, y=1cm]
  \path[use as bounding box] (-0.15,-0.10) rectangle (7.85,7.75);
  \coordinate (O) at (3.85,3.80);
  \def\R{2.75}
  \coordinate (A) at ($(O)+(220:\R)$);
  \coordinate (B) at ($(O)+(320:\R)$);
  \coordinate (P) at ($(O)+(90:\R)$);
  \fill[eleveazulclaro] (O) circle (\R);
  \draw[eleve aresta, line width=1.7pt] (O) circle (\R);
  \draw[eleve destaque, line width=3pt]
    ($(O)+(220:\R)$) arc[start angle=220,end angle=320,radius=\R];
  \draw[eleve aresta, line width=1.5pt] (O) -- (A) (O) -- (B);
  \draw[eleve aresta, line width=1.5pt] (P) -- (A) (P) -- (B);
  \draw[eleve destaque, line width=1.4pt]
    ($(O)+(220:0.72)$) arc[start angle=220,end angle=320,radius=.72];
  \draw[eleve destaque, line width=1.4pt]
    ($(P)+(238:.68)$) arc[start angle=238,end angle=302,radius=.68];
  \node[font=\Large\bfseries, text=elevelaranja] at (3.85,2.88) {$2\alpha$};
  \node[font=\Large\bfseries, text=elevelaranja] at (3.85,5.66) {$\alpha$};
  \EleveVertice{A}{below left}{A}
  \EleveVertice{B}{below right}{B}
  \EleveVertice{P}{above}{P}
  \EleveVertice{O}{above right}{O}
  \node[font=\large\bfseries, text=elevecinza] at (3.85,0.48)
    {mesmo arco $AB$};
\end{tikzpicture}

% FIGURA 7 — fig-07-angulo-de-segmento
\begin{tikzpicture}[eleve figura, x=1cm, y=1cm]
  \path[use as bounding box] (-0.15,-0.10) rectangle (8.15,7.95);
  \coordinate (O) at (3.60,3.80);
  \def\R{2.45}
  \coordinate (T) at ($(O)+(0:\R)$);
  \coordinate (A) at ($(O)+(132:\R)$);
  \fill[eleveazulclaro] (O) circle (\R);
  \draw[eleve aresta, line width=1.7pt] (O) circle (\R);
  \draw[eleve destaque, <->, line width=1.7pt] (T |- 0,0.35) -- (T |- 0,7.25);
  \draw[eleve aresta, line width=1.6pt] (T) -- (A);
  \draw[eleve auxiliar] (O) -- (T);
  \draw[elevecinza, line width=1pt] ($(T)+(-.28,0)$) -- ($(T)+(-.28,.28)$) -- ($(T)+(0,.28)$);
  \draw[eleve destaque, line width=2.8pt]
    ($(O)+(0:\R)$) arc[start angle=0,end angle=132,radius=\R];
  \draw[eleve destaque, line width=1.4pt]
    ($(T)+(90:.68)$) arc[start angle=90,end angle=156,radius=.68];
  \node[font=\Large\bfseries, text=elevelaranja] at (5.44,4.76) {$\alpha$};
  \node[font=\large\bfseries, text=elevecinza, rotate=90] at (6.45,3.80)
    {tangente};
  \node[font=\large\bfseries, text=eleveazul, rotate=-24] at (3.35,4.34)
    {corda};
  \node[font=\large\bfseries, text=elevelaranja] at (4.40,6.80)
    {arco correspondente};
  \EleveVertice{A}{above left}{A}
  \EleveVertice{T}{right}{T}
  \EleveVertice{O}{below left}{O}
\end{tikzpicture}

% FIGURA 8 — fig-08-angulo-excentrico-interior
\begin{tikzpicture}[eleve figura, x=1cm, y=1cm]
  \path[use as bounding box] (-0.15,-0.10) rectangle (8.05,7.95);
  \coordinate (O) at (3.95,3.90);
  \def\R{2.70}
  \coordinate (A) at ($(O)+(142:\R)$);
  \coordinate (B) at ($(O)+(322:\R)$);
  \coordinate (C) at ($(O)+(218:\R)$);
  \coordinate (D) at ($(O)+(38:\R)$);
  \coordinate (P) at (O);
  \fill[eleveazulclaro] (O) circle (\R);
  \draw[eleve aresta, line width=1.7pt] (O) circle (\R);
  \draw[eleve aresta, line width=1.6pt] (A) -- (B);
  \draw[eleve destaque, line width=1.6pt] (C) -- (D);
  \draw[eleve destaque, line width=2.8pt]
    ($(O)+(38:\R)$) arc[start angle=38,end angle=142,radius=\R];
  \draw[elevecinza, dashed, line width=2.3pt]
    ($(O)+(218:\R)$) arc[start angle=218,end angle=322,radius=\R];
  \draw[eleve destaque, line width=1.4pt]
    ($(P)+(38:.72)$) arc[start angle=38,end angle=142,radius=.72];
  \fill[elevecinza] (P) circle (2.5pt)
    node[font=\large\bfseries, text=elevecinza, below=7pt] {$P$};
  \node[font=\Large\bfseries, text=elevelaranja] at (3.95,4.92) {$\gamma_{int}$};
  \node[font=\large\bfseries, text=elevelaranja] at (3.95,7.18) {arco $\alpha_1$};
  \node[font=\large\bfseries, text=elevecinza] at (3.95,0.60) {arco $\alpha_2$};
  \EleveVertice{A}{above left}{A}
  \EleveVertice{B}{below right}{B}
  \EleveVertice{C}{below left}{C}
  \EleveVertice{D}{above right}{D}
\end{tikzpicture}

% FIGURA 9 — fig-09-angulo-excentrico-exterior
\begin{tikzpicture}[eleve figura, x=1cm, y=1cm]
  \path[use as bounding box] (-0.15,-0.10) rectangle (9.25,7.80);
  \coordinate (O) at (5.55,3.75);
  \def\R{2.40}
  \coordinate (P) at (0.55,3.75);
  \fill[eleveazulclaro] (O) circle (\R);
  \draw[eleve aresta, line width=1.7pt] (O) circle (\R);
  \draw[eleve aresta, ->, line width=1.6pt] (P) -- (8.55,7.35);
  \draw[eleve destaque, ->, line width=1.6pt] (P) -- (8.55,0.15);
  % Os arcos interceptados ficam em lados opostos da circunferência.
  \draw[eleve destaque, line width=2.8pt]
    ($(O)+(316:\R)$) arc[start angle=316,end angle=404,radius=\R];
  \draw[elevecinza, dashed, line width=2.3pt]
    ($(O)+(136:\R)$) arc[start angle=136,end angle=224,radius=\R];
  \draw[eleve destaque, line width=1.4pt]
    ($(P)+(-24:.92)$) arc[start angle=-24,end angle=24,radius=.92];
  \fill[elevelaranja] (P) circle (2.7pt)
    node[font=\large\bfseries, text=elevecinza, left=5pt] {$P$};
  \node[font=\LARGE\bfseries, text=elevelaranja] at (2.25,3.75) {$\gamma_{ext}$};
  \node[font=\large\bfseries, text=elevelaranja, align=left] at (8.00,3.75)
    {arco\\maior};
  \node[font=\large\bfseries, text=elevecinza, align=left] at (4.05,3.75)
    {arco\\menor};
\end{tikzpicture}

% FIGURA 10 — fig-10-potencia-entre-cordas
\begin{tikzpicture}[eleve figura, x=1cm, y=1cm]
  \path[use as bounding box] (-0.15,-0.10) rectangle (8.15,7.75);
  \coordinate (O) at (3.95,3.80);
  \def\R{2.70}
  \coordinate (A) at ($(O)+(165:\R)$);
  \coordinate (B) at ($(O)+(345:\R)$);
  \coordinate (C) at ($(O)+(230:\R)$);
  \coordinate (D) at ($(O)+(50:\R)$);
  \coordinate (P) at (O);
  \fill[eleveazulclaro] (O) circle (\R);
  \draw[eleve aresta, line width=1.7pt] (O) circle (\R);
  \draw[eleve aresta, line width=1.8pt] (A) -- (P) -- (B);
  \draw[eleve destaque, line width=1.8pt] (C) -- (P) -- (D);
  \fill[elevecinza] (P) circle (2.7pt)
    node[font=\large\bfseries, text=elevecinza, below right=4pt] {$P$};
  \node[font=\large\bfseries, text=eleveazul, rotate=-15] at ($(A)!0.5!(P)+(0,.47)$) {$PA$};
  \node[font=\large\bfseries, text=eleveazul, rotate=-15] at ($(P)!0.5!(B)+(0,.47)$) {$PB$};
  \node[font=\large\bfseries, text=elevelaranja, rotate=50] at ($(C)!0.5!(P)+(-.50,0)$) {$PC$};
  \node[font=\large\bfseries, text=elevelaranja, rotate=50] at ($(P)!0.5!(D)+(-.50,0)$) {$PD$};
  \EleveVertice{A}{left}{A}
  \EleveVertice{B}{right}{B}
  \EleveVertice{C}{below left}{C}
  \EleveVertice{D}{above right}{D}
\end{tikzpicture}

% FIGURA 11 — fig-11-potencia-entre-secantes
\begin{tikzpicture}[eleve figura, x=1cm, y=1cm]
  \path[use as bounding box] (-0.15,-0.10) rectangle (9.45,7.75);
  \coordinate (O) at (5.65,3.75);
  \def\R{2.45}
  \coordinate (P) at (0.55,3.75);
  \fill[eleveazulclaro] (O) circle (\R);
  \draw[eleve aresta, line width=1.7pt] (O) circle (\R);
  \draw[eleve aresta, line width=1.8pt] (P) -- (8.70,7.10);
  \draw[eleve destaque, line width=1.8pt] (P) -- (8.70,0.40);
  \coordinate (A) at (3.55,4.98);
  \coordinate (B) at (7.92,6.78);
  \coordinate (C) at (3.55,2.52);
  \coordinate (D) at (7.92,0.72);
  \EleveVertice{A}{above left}{A}
  \EleveVertice{B}{above right}{B}
  \EleveVertice{C}{below left}{C}
  \EleveVertice{D}{below right}{D}
  \fill[elevecinza] (P) circle (2.7pt)
    node[font=\large\bfseries, text=elevecinza, left=5pt] {$P$};
  \node[font=\Large\bfseries, text=eleveazul, rotate=22] at (2.20,5.15) {$PA$};
  \node[font=\Large\bfseries, text=eleveazul, rotate=22] at (6.15,6.80) {$PB$};
  \node[font=\Large\bfseries, text=elevelaranja, rotate=-22] at (2.20,2.32) {$PC$};
  \node[font=\Large\bfseries, text=elevelaranja, rotate=-22] at (6.15,0.70) {$PD$};
\end{tikzpicture}

% FIGURA 12 — fig-12-potencia-secante-tangente
\begin{tikzpicture}[eleve figura, x=1cm, y=1cm]
  \path[use as bounding box] (-0.15,-0.10) rectangle (8.65,7.35);
  \coordinate (O) at (4.65,3.15);
  \def\R{2.00}
  \coordinate (P) at (0.65,3.15);
  \coordinate (A) at (2.65,3.15);
  \coordinate (B) at (6.65,3.15);
  \coordinate (T) at (3.65,4.882);
  \fill[eleveazulclaro] (O) circle (\R);
  \draw[eleve aresta, line width=1.7pt] (O) circle (\R);
  \draw[eleve aresta, ->, line width=1.8pt] (P) -- (7.75,3.15);
  \draw[eleve destaque, line width=1.9pt] (P) -- (T);
  \draw[eleve auxiliar] (O) -- (T);
  \draw[elevecinza, line width=1pt]
    ($(T)!0.16!(O)$) -- ($(T)!0.16!(O)+(150:.28)$) -- ($(T)+(150:.28)$);
  \EleveVertice{A}{below}{A}
  \EleveVertice{B}{below}{B}
  \EleveVertice{T}{above right}{T}
  \fill[elevecinza] (P) circle (2.7pt)
    node[font=\large\bfseries, text=elevecinza, left=5pt] {$P$};
  \EleveVertice{O}{above right}{O}
  \node[font=\large\bfseries, text=eleveazul] at (1.67,2.64) {$PA$};
  \node[font=\large\bfseries, text=eleveazul] at (5.56,2.64) {$PB$};
  \node[font=\Large\bfseries, text=elevelaranja, rotate=24] at (2.10,4.78) {$PT$};
  \node[font=\large\bfseries, text=elevelaranja, align=left] at (6.85,5.55)
    {tangente\\em $T$};
\end{tikzpicture}

\end{document}
